\subsection{Модуль Ledbar}

Для управления светодиодами в системе Geoskan Pioneer используется встроенный модуль \texttt{Ledbar}. Он предоставляет простой интерфейс (API) для работы с цветом.

\begin{codebox}[Основные функции]
\begin{itemize}
  \item \texttt{Ledbar.new(N)} — конструктор. Создаёт объект для управления лентой из \texttt{N} светодиодов.
  \begin{itemize}
      \item \textit{Пример:} \texttt{local leds = Ledbar.new(29)}
  \end{itemize}
  
  \item \texttt{leds:set(i, r, g, b)} — устанавливает цвет конкретного светодиода.
  \begin{itemize}
      \item \texttt{i} — индекс светодиода (от \texttt{0} до \texttt{N-1}).
      \item \texttt{r, g, b} — яркость каналов в диапазоне от \texttt{0.0} (выключен) до \texttt{1.0} (максимум).
  \end{itemize}
\end{itemize}
\end{codebox}

\begin{warnbox}[Индексация]
Обратите внимание: нумерация светодиодов начинается с \textbf{нуля}. Если вы создали ленту на 29 элементов, допустимые индексы — от 0 до 28. Обращение к индексу вне этого диапазона (например, к индексу 29) \textbf{не вызовет ошибку}: вызов \texttt{leds:set(...)} будет молча проигнорирован, и светодиод просто не загорится. Это коварнее явной ошибки — программа продолжит работать как ни в чём не бывало, и вы не получите никакой подсказки. Поэтому если светодиод не светится, первым делом проверьте индекс.
\end{warnbox}

Также часто используется вспомогательная функция для заливки всей ленты одним цветом:
\begin{codebox}[Вспомогательная функция]
\begin{codefile}{lua/examples/05_leds/leds_api_change_color.lua}
\end{codefile}
\end{codebox}

\textbf{Пояснение.} Функция \verb|changeColor(color)| принимает таблицу \verb|{r, g, b}| и в цикле вызывает \verb|leds:set(i, ...)| для всех индексов \verb|0..N-1|. Основной принцип — вынести повторяющийся код установки одного и того же цвета в отдельную функцию, чтобы проще переиспользовать его в разных сценариях.
